\subsection{Knowledge graph under data scarcity}
\label{sec:kgeff}
Because ontology-based representations are expected to help most when
training data are limited, we retrained TKGN with and without the knowledge
graph on 5, 10 and 25\% of the training patients in three grouped repeats,
with identical patient subsets and seeds for both variants
(Table~\ref{tab:kg}). The hypothesis was not supported. For readmission the
paired differences were close to zero at every size (\KgReadmitFive{},
\KgReadmitTen{} and \KgReadmitTwentyFive{} AUROC at 5, 10 and 25\%), with
standard deviations several times larger than the means. For escalation the
flat-embedding variant was never worse on average (\KgEscFive{},
\KgEscTen{} and \KgEscTwentyFive{}). On this cohort, whose diagnosis codes
are mostly truncated to three characters, sharing information along the
ICD-9-CM hierarchy and the co-morbidity graph therefore did not improve
discrimination at any training-set size we examined.

\begin{table}[t]
  \centering
  \caption{TKGN with and without the knowledge graph (KG) trained on subsets
  of the training patients. Mean test AUROC over three grouped repeats with
  paired subsets and seeds; $\Delta$ = TKGN minus TKGN without KG; Wins:
  repeats in which the KG variant was better.}
  \label{tab:kg}
  \scriptsize
  \setlength{\tabcolsep}{3pt}
  \begin{tabular}{@{}llcccc@{}}
\toprule
Task & Training share & TKGN & TKGN $-$ KG & $\Delta$AUROC (mean$\pm$SD) & Wins \\
\midrule
30-day readmission & 5\% & 0.6387 & 0.6395 & -0.0009$\pm$0.0063 & 1/3 \\
 & 10\% & 0.6495 & 0.6487 & +0.0008$\pm$0.0054 & 2/3 \\
 & 25\% & 0.6717 & 0.6715 & +0.0003$\pm$0.0049 & 2/3 \\
 & 100\% & 0.6847 & 0.6857 & -0.0011$\pm$0.0017 & 0/3 \\
\midrule
Treatment escalation & 5\% & 0.6049 & 0.6201 & -0.0152$\pm$0.0205 & 1/3 \\
 & 10\% & 0.6331 & 0.6336 & -0.0005$\pm$0.0031 & 1/3 \\
 & 25\% & 0.6585 & 0.6596 & -0.0012$\pm$0.0030 & 1/3 \\
 & 100\% & 0.6781 & 0.6819 & -0.0038$\pm$0.0021 & 0/3 \\
\bottomrule
\end{tabular}

\end{table}
